\documentclass[10pt]{article}
\usepackage[T1]{fontenc}
\usepackage[utf8]{inputenc}
\usepackage[french]{babel}
\usepackage{amsmath,amssymb,mathrsfs,mathpazo}
\usepackage{geometry}
\geometry{a4paper,margin=2cm}
\usepackage{graphicx,array,tabularx,tabularray,colortbl,multirow,multicol,booktabs,enumitem,pifont,ulem,fancybox,diagbox,cancel}
\usepackage{pstricks,pstricks-add,pst-plot,pst-tree,pst-node,pst-text,pst-func,pst-math,pst-bspline,pst-3dplot,pst-eucl,pst-solides3d}
\let\pscalculate\undefined
\usepackage{pst-bezier}
\usepackage{tikz,pgfplots,tkz-tab}
\usetikzlibrary{arrows,arrows.meta,patterns,calc,positioning,trees,decorations.markings,shapes.geometric,shapes.misc,shapes.arrows,intersections,decorations.pathmorphing,decorations.pathreplacing,backgrounds,shadows,babel}
\pgfplotsset{compat=1.18}
\usepgfplotslibrary{fillbetween}
\usepackage[np]{numprint}
\usepackage{listings,makecell,pas-tableur,twoopt}
\let\barre\undefined
\usepackage[pstricks]{tablvar}
\definecolor{darkgreen}{RGB}{0,100,0}
\providecommand{\red}{\color{red}}
\providecommand{\blue}{\color{blue}}
\providecommand{\green}{\color{green}}
\providecommand{\black}{\color{black}}
\providecommand{\white}{\color{white}}
\providecommand{\gray}{\color{gray}}
\providecommand{\cyan}{\color{cyan}}
\providecommand{\magenta}{\color{magenta}}
\providecommand{\yellow}{\color{yellow}}
\providecommand{\violet}{\color{violet}}
\providecommand{\orange}{\color{orange}}
\providecommand{\euro}{\texteuro}
\providecommand{\EUR}{\texteuro}
\providecommand{\ieme}{\textsuperscript{e}}
\providecommand{\degres}{\ensuremath{^\circ}}
\providecommand{\degree}{\ensuremath{^\circ}}
\providecommand{\vv}[1]{\overrightarrow{#1}}
\providecommand{\Coord}[1]{#1}
\providecommand{\decosix}{}
\providecommand{\starredbullet}{\ensuremath{\bullet}}
\pagestyle{empty}
\setlength{\parindent}{0pt}
\usepackage{adjustbox}\DeclareUnicodeCharacter{2236}{:}\DeclareUnicodeCharacter{2190}{\ensuremath{\leftarrow}}
\providecommand{\R}{\mathbb{R}}
\providecommand{\N}{\mathbb{N}}
\providecommand{\Z}{\mathbb{Z}}
\providecommand{\Q}{\mathbb{Q}}
\providecommand{\D}{\mathbb{D}}
\providecommand{\C}{\mathbb{C}}
\providecommand{\e}{{\mathrm{e}}}
\providecommand{\vect}[1]{\overrightarrow{#1}}
\providecommand{\vectt}[1]{\overrightarrow{\mathrm{#1}}}
\providecommand{\vvt}[1]{\overrightarrow{\mathrm{#1}}}
\providecommand{\np}[2][]{{#2~#1}}
\providecommand{\vv}[1]{{\overrightarrow{#1}}}
\renewcommand{\d}{{\mathrm{d}}}
\providecommand{\degres}{{^\circ}}
\providecommand{\degre}{{^\circ}}
\providecommand{\blue}{\color{blue}}
\providecommand{\red}{\color{red}}
\providecommand{\enskip}{\quad}
\providecommand{\llbracket}{\lbrack\!\lbrack}
\providecommand{\rrbracket}{\rbrack\!\rbrack}
\providecommand{\euro}{{\text{€}}}
\providecommand{\no}{{\text{n°}}}
\providecommand{\interval}[2]{{[#1;#2]}}
\providecommand{\intervFO}[2]{{[#1;#2[}}
\providecommand{\dotfill}{\ldots}

\providecommand{\ds}{\displaystyle}
\providecommand{\textding}[1]{#1}
\providecommand{\Oij}{\(\mathrm{O};\vec{\imath},\vec{\jmath}\)}
\providecommand{\Oijk}{\(\mathrm{O};\vec{\imath},\vec{\jmath},\vec{k}\)}
\providecommand{\Ouv}{\(\mathrm{O};\vec{u},\vec{v}\)}
\providecommand{\intervFF}[2]{\ensuremath{[#1;#2]}}
\renewcommand{\euro}{\ifmmode\text{\texteuro}\else\texteuro\fi}
\DeclareUnicodeCharacter{00B0}{\ifmmode^{\circ}\else\textdegree\fi}
\DeclareUnicodeCharacter{20AC}{\ifmmode\text{\texteuro}\else\texteuro\fi}
% Énoncés : Éducation nationale. Transcription et corrigés : APMEP, Denis Vergès et contributeurs. Découpage et classement : CoopMaths, Sébastien Lozano, Éric Elter et contributeurs. Adaptation de la présentation : Maths Explorer.
% CC BY-SA : https://coopmaths.fr/www/about/
% https://coopmaths.fr/alea/static/eam/2026/tex/eam_technologique_2026_09_antilles_2.tex
\begin{document}
\noindent\textit{APMEP / CoopMaths. CC BY-SA. Présentation adaptée par Maths Explorer.}\par\medskip
\textbf{Pour chaque affirmation, répondre par VRAI ou FAUX, en justifiant la réponse.}

\medskip

\begin{minipage}[t]{0.58\linewidth}
	Au 1\up{er} janvier 2025, un étang contient \np{3000}~m$^3$ d'eau.

	Chaque année, l'étang perd 10\,\% de son volume d'eau.

	On modélise l'évolution du volume d'eau de cet étang par une suite $\left(u_n\right)$ où, pour tout entier naturel $n$, $u_n$ désigne la quantité d'eau, en mètre cube, contenue dans l'étang, le 1\up{er} janvier de l'année $2025+n$.

	Ainsi, on a $u_0 = \np{3000}$.

	La feuille de calcul ci-contre donne les valeurs des premiers termes de cette suite $\left(u_n\right)$.
\end{minipage}
\hfill
\begin{minipage}[t]{0.38\linewidth}
	\vspace{0pt}
	\begin{tikzpicture}[scale=0.62,transform shape]
		\tableur*[15]{A/2.5cm,B/2.5cm}
		\celtxt*[align=center]{A}{1}{$n$} \celtxt*[align=center]{B}{1}{$u_n$}
		\celtxt*[align=center]{A}{2}{$0$} \celtxt*[align=center]{B}{2}{$3000$}
		\celtxt*[align=center]{A}{3}{$1$} \celtxt*[align=center]{B}{3}{$2700$}
		\celtxt*[align=center]{A}{4}{$2$} \celtxt*[align=center]{B}{4}{$2430$}
		\celtxt*[align=center]{A}{5}{$3$} \celtxt*[align=center]{B}{5}{$2187$}
		\celtxt*[align=center]{A}{6}{$4$} \celtxt*[align=center]{B}{6}{$1968{,}3$}
		\celtxt*[align=center]{A}{7}{$5$} \celtxt*[align=center]{B}{7}{$1771{,}47$}
		\celtxt*[align=center]{A}{8}{$6$} \celtxt*[align=center]{B}{8}{$1594{,}323$}
		\celtxt*[align=center]{A}{9}{$7$} \celtxt*[align=center]{B}{9}{$1434{,}8907$}
		\celtxt*[align=center]{A}{10}{$8$} \celtxt*[align=center]{B}{10}{$1291{,}40163$}
		\celtxt*[align=center]{A}{11}{$9$} \celtxt*[align=center]{B}{11}{$1162{,}26147$}
		\celtxt*[align=center]{A}{12}{$10$} \celtxt*[align=center]{B}{12}{$1046{,}03532$}
		\celtxt*[align=center]{A}{13}{$11$} \celtxt*[align=center]{B}{13}{$941{,}431788$}
		\celtxt*[align=center]{A}{14}{$12$} \celtxt*[align=center]{B}{14}{$847{,}288609$}
		\celtxt*[align=center]{A}{15}{$13$} \celtxt*[align=center]{B}{15}{$762{,}559748$}
	\end{tikzpicture}
\end{minipage}

\begin{enumerate}
	\item \textbf{Affirmation}

	Pour obtenir les termes de la suite $\left(u_n\right)$, la formule saisie dans la cellule \textsf{B3} puis étirée vers le bas est : \fbox{\small\textsf{= 0,90 * A3}}

	\item \textbf{Affirmation}

	La suite $\left(u_n\right)$ est une suite arithmétique de raison $r = 0,90$ et de 1\up{er} terme $u_0 = \np{3000}$.
\end{enumerate}

\medskip

\begin{minipage}[t]{0.58\linewidth}
	Au 1\up{er} janvier 2025, un second étang contient quant à lui \np{4000}~m$^3$ d'eau. Chaque année, cet étang perd 20\,\% de son volume d'eau initial mais il est naturellement alimenté par un cours d'eau qui augmente son volume d'eau de 140~m$^3$ par an.

	On peut également modéliser l'évolution du volume d'eau de cet étang par une suite $\left(v_n\right)$ où, pour tout entier naturel $n$, $v_n$ désigne la quantité d'eau, en mètre cube, contenue dans l'étang, le 1\up{er} janvier de l'année $2025+n$.

	Ainsi, on a $v_0 = \np{4000}$.

	La feuille de calcul ci-contre donne les valeurs des premiers termes de cette suite $\left(v_n\right)$.
\end{minipage}
\hfill
\begin{minipage}[t]{0.38\linewidth}
	\vspace{0pt}
	\begin{tikzpicture}[scale=0.62,transform shape]
		\tableur*[15]{A/2.5cm,B/2.5cm}
		\fill[darkgray] (cellB-3.south west) rectangle (cellB-3.north east);
		\celtxt*[align=center]{A}{1}{$n$} \celtxt*[align=center]{B}{1}{$v_n$}
		\celtxt*[align=center]{A}{2}{$0$} \celtxt*[align=center]{B}{2}{$4000$}
		\celtxt*[align=center]{A}{3}{$1$}
		\celtxt*[align=center]{A}{4}{$2$} \celtxt*[align=center]{B}{4}{$2812$}
		\celtxt*[align=center]{A}{5}{$3$} \celtxt*[align=center]{B}{5}{$2389{,}6$}
		\celtxt*[align=center]{A}{6}{$4$} \celtxt*[align=center]{B}{6}{$2051{,}68$}
		\celtxt*[align=center]{A}{7}{$5$} \celtxt*[align=center]{B}{7}{$1781{,}344$}
		\celtxt*[align=center]{A}{8}{$6$} \celtxt*[align=center]{B}{8}{$1565{,}0752$}
		\celtxt*[align=center]{A}{9}{$7$} \celtxt*[align=center]{B}{9}{$1392{,}06016$}
		\celtxt*[align=center]{A}{10}{$8$} \celtxt*[align=center]{B}{10}{$1253{,}64813$}
		\celtxt*[align=center]{A}{11}{$9$} \celtxt*[align=center]{B}{11}{$1142{,}9185$}
		\celtxt*[align=center]{A}{12}{$10$} \celtxt*[align=center]{B}{12}{$1054{,}3348$}
		\celtxt*[align=center]{A}{13}{$11$} \celtxt*[align=center]{B}{13}{$983{,}467842$}
		\celtxt*[align=center]{A}{14}{$12$} \celtxt*[align=center]{B}{14}{$926{,}774273$}
		\celtxt*[align=center]{A}{15}{$13$} \celtxt*[align=center]{B}{15}{$881{,}419419$}
	\end{tikzpicture}
\end{minipage}

\begin{enumerate}[start=3]
	\item \textbf{Affirmation}

	On a $v_1 = \np{3340}$.

	\item \textbf{Affirmation}

	Les volumes des deux étangs passeront simultanément en dessous de \np{1000}~m$^3$ à partir de 2036.
\end{enumerate}

\medskip

\end{document}
